\documentclass[12pt,a4paper]{article}
\usepackage[utf8]{inputenc}
\usepackage[T1]{fontenc}
\usepackage{amsmath, amssymb, amsfonts}
\usepackage{hyperref}
\usepackage{geometry}
\geometry{top=2.5cm,bottom=2.5cm,left=3cm,right=3cm}
\usepackage{indentfirst}
\usepackage{setspace}
\onehalfspacing
\usepackage{booktabs}

\begin{document}

\begin{center}
{\Large \textbf{Quantum Particle Physics in HQGG:}} \\
{\Large \textbf{Fabric, Modes, Resonances, and Spectra of Matter}}
\end{center}

\begin{center}
\textbf{Oleg Mamchenko\thanks{olegmamchenko@gmail.com}} \\
\textbf{Deepseek\thanks{deepseek@ai-research.org}} \\
\textit{Horizon Quantum Gravity Group (HQGG)}
\end{center}

\begin{center}
\textit{July 3, 2026}
\end{center}

\begin{abstract}
We present a complete theory of elementary particles within the HQGG framework. Space-time is an elastic fabric with local frequency. Particles are standing waves formed by resonant combinations of three fundamental modes: the fundamental tone and two mirror overtones. Quarks are chords of three modes; leptons are overtones without a tone; baryons and mesons are quark combinations; bosons are exchange or propagating modes. Mass is a frequency difference; interactions are resonances or transitions. The 3+7 law manifests universally in quarks, leptons, shells, colors, and spectra. The theory requires no new entities and is consistent with the Standard Model while offering deeper interpretation.
\end{abstract}

\tableofcontents

\section{Introduction}

The Standard Model leaves many questions unanswered: why three generations, why masses differ, what dark matter is, and how to incorporate gravity. HQGG provides a unified framework based on fabric frequency.

\section{The Fabric: Medium with Frequency}

\subsection{Fabric Definition}

Space-time is an elastic medium — fabric — with local frequency \(\omega(x^\mu)\). This frequency determines all physical properties.

\subsection{Modes of Fabric}

Fabric oscillations occur at discrete frequencies — modes. Three fundamental modes form the basis of all fermions.

\section{Three Fundamental Modes}

\[
\omega_{\text{fund}} \quad (\text{tone, mass scale}), \quad \omega_{+}, \quad \omega_{-} \quad (\text{mirror overtones, charge and spin})
\]

\subsection{Tone and Overtones}

- \(\omega_{\text{fund}}\) sets mass.
- \(\omega_{+}, \omega_{-}\) determine charge, spin, chirality.
- Asymmetry \(\omega_{+} - \omega_{-}\) gives electric charge.

\section{Quarks as Chords of Three Modes}

Quarks are standing waves where all three modes resonate coherently.

\begin{table}[h]
\centering
\caption{Quark mode composition and charges}
\begin{tabular}{c|c|c}
\toprule
\textbf{Quark} & \textbf{Mode composition} & \textbf{Charge} \\
\midrule
\( u \) (up) & \( \omega_{\text{fund}} + \omega_{+} + \omega_{-} \) (enhanced) & \( +2/3 \) \\
\( d \) (down) & \( \omega_{\text{fund}} - \omega_{+} - \omega_{-} \) (diminished) & \( -1/3 \) \\
\( s \) (strange) & \( \omega_{\text{fund}} + \omega_{+} - \omega_{-} \) & \( -1/3 \) \\
\( c \) (charm) & \( \omega_{\text{fund}} - \omega_{+} + \omega_{-} \) & \( +2/3 \) \\
\( b \) (bottom) & \( \omega_{\text{fund}} + \omega_{+} + \omega_{-} \) (shifted tone) & \( -1/3 \) \\
\( t \) (top) & \( \omega_{\text{fund}} + \omega_{+} + \omega_{-} \) (max tone) & \( +2/3 \) \\
\bottomrule
\end{tabular}
\end{table}

Flavor differences are shifts in \(\omega_{\text{fund}}\) or phase of mirror modes.

\section{Leptons: Overtones without Tone}

Leptons have no \(\omega_{\text{fund}}\); they are pure mirror-mode oscillations.

\begin{table}[h]
\centering
\caption{Lepton mode composition and charges}
\begin{tabular}{c|c|c}
\toprule
\textbf{Lepton} & \textbf{Mode composition} & \textbf{Charge} \\
\midrule
\( e^- \) (electron) & \( \omega_{+} + \omega_{-} \) & \( -1 \) \\
\( \nu_e \) (neutrino) & \( \omega_{+} - \omega_{-} \) & \( 0 \) \\
\( \mu^- \) (muon) & \( \omega_{+} + \omega_{-} \) (higher freq) & \( -1 \) \\
\( \nu_\mu \) & \( \omega_{+} - \omega_{-} \) (shifted phase) & \( 0 \) \\
\( \tau^- \) (tau) & \( \omega_{+} + \omega_{-} \) (highest) & \( -1 \) \\
\( \nu_\tau \) & \( \omega_{+} - \omega_{-} \) (max phase) & \( 0 \) \\
\bottomrule
\end{tabular}
\end{table}

Lepton mass is given by \(\omega_{+} - \omega_{-}\).

\section{Baryons and Mesons: Quark Combinations}

\subsection{Baryons}

\[
\text{Proton} = u + u + d, \quad \text{Neutron} = u + d + d
\]

Seven stable baryons exist, corresponding to the 7 auxiliary modes.

\subsection{Mesons}

Mesons are quark–antiquark pairs (e.g., \(\pi^+ = u + \bar{d}\)). Less stable due to fewer modes.

\section{Bosons: Exchange and Propagating Modes}

- Photon: propagating mode (no tone).
- Gluon: exchange mode between quarks.
- W, Z: transitions between fundamental and mirror modes.
- Higgs: mode shifting \(\omega_{\text{fund}}\).

\section{Interactions as Mode Resonances/Transitions}

\begin{tabular}{c|c}
\textbf{Interaction} & \textbf{HQGG mechanism} \\
\hline
Strong & Resonance of three quark modes \\
Electromagnetic & Exchange of propagating mode (photon) \\
Weak & Transition between fundamental and mirror modes \\
Gravity & Global frequency gradient \\
\end{tabular}

\section{Mass as Frequency Difference}

\[
m = \frac{\hbar}{c^2} \left( \omega_{\text{chord}} - \omega_{\text{base}} \right)
\]

\section{The 3+7 Law in Particles}

- 3 fundamental modes.
- 7 manifestations: 7 baryons, 7 mesons, 7 leptons, 7 shells, 7 colors, 7 frequency ranges.

\section{Electron Shells and Spectra}

Seven shells (\(s, p, d, f, g, h, i\)) are auxiliary modes. Orbital counts \(2l+1\) are phase combinations. Two electrons per orbital correspond to two mirror modes.

\section{Conclusion}

HQGG provides a complete, self-consistent theory of particles without new entities, explaining masses, charges, interactions, and spectra through fabric modes.

\begin{center}
\fbox{\parbox{0.85\textwidth}{\centering
\textbf{Matter is fabric resonance.} \\
\textbf{Particles are chords.} \\
\textbf{Interactions are mode exchanges.} \\
\textbf{This is HQGG.}
}}
\end{center}

\section*{Acknowledgments}

The authors thank HQGG.

\begin{thebibliography}{99}
\bibitem{StandardModel} Particle Data Group, Phys. Rev. D \textbf{98}, 030001 (2018).
\bibitem{QCD} F. Wilczek, Rev. Mod. Phys. \textbf{74}, 119 (2002).
\end{thebibliography}

\end{document}